%% =====================================================================
%%  Plan du tome 1 — Premiers pas (A1) : 5 parties, 20 unités, 5 bilans
%%  \planPartie{n}{titre}{titre anglais}{couleur}{première}{dernière}{présentation}
%%  \planUnite{n}{titre anglais}{sous-titre}{communication}{grammaire}
%%            {vocabulaire}{prononciation}{culture}
%% =====================================================================
\planPartie{1}{Premiers contacts}{First steps}{otNavy}{1}{4}
  {Saluer, se présenter, dire d'où l'on vient, épeler son nom, donner son numéro et présenter
   sa famille : les toutes premières phrases pour entrer en contact.}
\planUnite{1}{Hello!}{Saluer et se présenter}
  {Saluer, prendre congé, se présenter, demander des nouvelles}
  {Les pronoms sujets ; \en{be} au présent (affirmatif)}
  {Salutations et politesse}
  {Le son \phon{h}}
  {\en{Hi}, \en{hello} ou \en{good morning} ? \en{Mr.}, \en{Mrs.}, \en{Ms.}}
\planUnite{2}{Where are you from?}{Dire d'où l'on vient}
  {Dire sa nationalité, sa ville, les langues que l'on parle}
  {\en{be} : négation, questions et réponses courtes}
  {Pays, nationalités, langues}
  {L'accent de mot}
  {Le monde anglophone}
\planUnite{3}{What's your number?}{Les lettres et les nombres}
  {Épeler son nom, donner son âge, son téléphone, son e-mail}
  {Les mots interrogatifs ; \en{a} et \en{an}}
  {L'alphabet, les nombres de 0 à 100}
  {Les lettres qui piègent : A, E, I, G, J}
  {Téléphones et adresses e-mail}
\planUnite{4}{This is my family}{Présenter sa famille}
  {Présenter ses proches, parler de sa famille}
  {Les adjectifs possessifs ; le génitif \en{'s}}
  {La famille, les prénoms}
  {Le son \phon{ð} de \en{this}}
  {Prénoms, surnoms et familles}
\planPartie{2}{La vie de tous les jours}{Everyday life}{otGreenDark}{5}{8}
  {Nommer les objets, raconter sa journée, dire ce qu'on aime et décrire son logement :
   l'anglais du quotidien, avec le présent simple.}
\planUnite{5}{What's this?}{Les objets autour de moi}
  {Nommer, montrer et décrire des objets}
  {\en{a}, \en{an}, \en{the} ; le pluriel ; \en{this}, \en{that}, \en{these}, \en{those}}
  {Objets du quotidien, couleurs}
  {Le pluriel : \phon{s}, \phon{z}, \phon{ɪz}}
  {Anglais américain et britannique}
\planUnite{6}{A day in my life}{Raconter sa journée}
  {Décrire sa routine, dire l'heure}
  {Le présent simple (affirmatif) ; le \en{-s} de la 3\textsuperscript{e} personne}
  {Activités quotidiennes, l'heure, les jours}
  {Le \en{-s} de la 3\textsuperscript{e} personne}
  {\en{a.m.} et \en{p.m.}, la journée de travail}
\planUnite{7}{Do you like coffee?}{Dire ce qu'on aime}
  {Exprimer ses goûts, poser des questions}
  {Présent simple : négation et questions (\en{do}, \en{does})}
  {Nourriture, boissons, loisirs}
  {Les formes faibles : \en{do you}, \en{does he}}
  {Petit-déjeuner et \en{coffee shops}}
\planUnite{8}{Home sweet home}{Décrire son logement}
  {Décrire un logement, situer des objets}
  {\en{There is}, \en{there are} ; les prépositions de lieu}
  {Les pièces, les meubles}
  {Le son \phon{θ} de \en{three}}
  {\en{Apartment} ou \en{flat} ? Se loger à New York}
\planPartie{3}{En ville}{Around town}{otOrangeDark}{9}{12}
  {Trouver son chemin, faire des achats, commander au café et parler de ses talents :
   se débrouiller seul en ville.}
\planUnite{9}{Excuse me, where's the bank?}{Demander son chemin}
  {Demander et indiquer un chemin}
  {L'impératif ; \en{next to}, \en{opposite}, \en{between}}
  {La ville, les lieux, les transports}
  {\phon{ɪ} ou \phon{iː} : \en{ship}, \en{sheep}}
  {Se repérer à Manhattan}
\planUnite{10}{How much is it?}{Faire des achats}
  {Demander un prix, acheter, payer}
  {\en{How much}, \en{how many} ; \en{some}, \en{any}}
  {Magasins, vêtements, prix}
  {\en{-teen} ou \en{-ty} ? \en{thirteen}, \en{thirty}}
  {Dollars, cents et taxes}
\planUnite{11}{Can I have a coffee, please?}{Au café et au restaurant}
  {Commander, demander poliment}
  {\en{Can I…?} ; \en{I'd like}, \en{Would you like…?}}
  {Le menu, les plats, les boissons}
  {Les formes faibles de \en{can}}
  {Le pourboire (\en{tipping})}
\planUnite{12}{I can swim!}{Dire ce qu'on sait faire}
  {Parler de ses capacités et de ses loisirs}
  {\en{can}, \en{can't} ; les adverbes (\en{well}, \en{very well})}
  {Sports, loisirs, talents}
  {\en{can} \phon{kæn} ou \en{can't} \phon{kænt} ?}
  {Le sport aux États-Unis}
\planPartie{4}{Maintenant et bientôt}{Now and next}{otViolet}{13}{16}
  {Décrire ce qui se passe, parler du temps qu'il fait, inviter et faire des projets :
   le présent en \en{-ing} et le futur proche.}
\planUnite{13}{What are you doing?}{Dire ce qui se passe}
  {Décrire une scène, parler au téléphone}
  {Le présent en \en{be + -ing} ; présent simple ou en \en{-ing} ?}
  {Activités, vêtements}
  {Le son \phon{ŋ} de \en{-ing}}
  {La politesse au téléphone}
\planUnite{14}{What's the weather like?}{Le temps et les saisons}
  {Parler de la météo, donner une date}
  {Les nombres ordinaux ; \en{in}, \en{on}, \en{at}}
  {Météo, mois, saisons}
  {Le son \phon{ɝː} de \en{first}}
  {Degrés Fahrenheit et jours fériés}
\planUnite{15}{Let's go out!}{Inviter et sortir}
  {Proposer, inviter, accepter, refuser poliment}
  {\en{Let's…}, \en{Would you like to…?} ; les pronoms compléments}
  {Sorties, spectacles, rendez-vous}
  {L'intonation pour accepter ou refuser}
  {Halloween et Thanksgiving}
\planUnite{16}{My plans}{Parler de ses projets}
  {Parler de ses projets et de ses intentions}
  {\en{be going to} ; les expressions du futur}
  {Voyages, vacances, projets}
  {\en{gonna}, \en{wanna} : l'anglais parlé}
  {Le \en{road trip} américain}
\planPartie{5}{Hier et aujourd'hui}{Then and now}{otTeal}{17}{20}
  {Dire où l'on était, raconter son week-end ou un voyage, et faire le bilan :
   le prétérit, pour parler du passé.}
\planUnite{17}{Where were you?}{Parler d'hier}
  {Dire où l'on était, décrire le passé}
  {\en{was}, \en{were} ; \en{there was}, \en{there were} ; \en{ago}}
  {Les expressions de temps du passé}
  {Les formes faibles de \en{was} et \en{were}}
  {Une brève histoire de New York}
\planUnite{18}{What did you do last weekend?}{Raconter son week-end}
  {Raconter des événements passés}
  {Le prétérit des verbes réguliers ; \en{did}, \en{didn't}}
  {Activités du week-end}
  {La terminaison \en{-ed} : \phon{t}, \phon{d}, \phon{ɪd}}
  {Le \en{brunch} du dimanche}
\planUnite{19}{A trip to Boston}{Raconter un voyage}
  {Raconter un voyage, poser des questions sur le passé}
  {Le prétérit des verbes irréguliers}
  {Voyage, transports, hébergement}
  {Les verbes irréguliers les plus courants}
  {Boston et la Nouvelle-Angleterre}
\planUnite{20}{Happy birthday, Léa!}{Faire le point}
  {Parler de soi au passé, au présent et au futur}
  {Révision : les temps et les questions}
  {Fêtes, cadeaux, souhaits}
  {Le rythme de la phrase anglaise}
  {Anniversaires et fêtes}
